\documentclass[11pt]{article}
\usepackage[margin=1in]{geometry}
\usepackage{amsmath,amssymb,amsthm}
\usepackage{hyperref}
\hypersetup{colorlinks=true,linkcolor=blue,citecolor=blue,urlcolor=blue,
  pdftitle={A two-dimensional complex with embedding-dimension gap two},
  pdfauthor={Alec Kriebel}}
\newtheorem{theorem}{Theorem}
\newtheorem{lemma}[theorem]{Lemma}
\newcommand{\R}{\mathbb R}
\newcommand{\PL}{\mathrm{PL}}
\newcommand{\lin}{\mathrm{lin}}
\newcommand{\sd}{\operatorname{sd}}
\newcommand{\conv}{\operatorname{conv}}
\title{A two-dimensional complex with\\embedding-dimension gap two}
\author{Alec Kriebel\\
  \small Independent researcher\\
  \small ORCID: \href{https://orcid.org/0009-0001-9320-500X}{0009-0001-9320-500X}}
\date{October 1, 2026 \quad (unrefereed research note, version 1.0)}
\begin{document}
\maketitle
\begin{abstract}
We prove that a finite two-dimensional simplicial complex can have minimum
piecewise-linear embedding dimension three and minimum original-simplex
linear embedding dimension five. This answers affirmatively the
minimum-dimension question posed by Brehm in 2006, already in dimension two.
The construction makes the random affine obstruction used by Newman
uniform under small facet deletions, then removes triangles sharing an
edge. The resulting linear hypergraph admits a PL embedding in three
dimensions. We include an elementary embedding of its first barycentric
subdivision and a finite parameter choice with exact arithmetic checks.
The existence proof does not generate an explicit facet list.
\end{abstract}

\section{The question and the result}
For a finite simplicial complex $K$, let $e_{\PL}(K)$ be the least
Euclidean dimension admitting a PL embedding of its polyhedron, and
$e_{\lin}(K)$ the least dimension admitting an embedding affine on every
\emph{original} simplex. Subdivision is permitted for the former invariant.
Brehm asked whether these minimum dimensions can differ by more than
one, with particular interest in dimensions two and three
\cite[p.~702]{Brehm06}.

\begin{theorem}\label{thm:main}
There exists a finite two-dimensional simplicial complex $K$ such that
\[
e_{\PL}(K)=3,\qquad e_{\lin}(K)=5.
\]
Moreover, its first barycentric subdivision embeds linearly in $\R^3$.
\end{theorem}

No manifold hypothesis is imposed in the question or the theorem.
Earlier Brehm--Sarkaria constructions separate the two embedding notions
in a fixed ambient dimension \cite{BS92}. Newman's random-complex
method gives such a separation in even ambient dimension
\cite{Newman23}; we use its balanced-Radon, order-type and probability
mechanism in $\R^4$. Lee--Nevo prove PL embeddability of linear uniform
hypergraphs \cite[Lemma~3.1]{LN}, with the relevant lemma already in
their July 2023 preprint. What the argument below adds to these checked
conclusions is a nonembedding guarantee surviving \emph{every} small
deletion, so an adaptive cleanup can be performed. Two independent
bounded literature audits found no earlier result subsuming the exact
conclusion; this is not a categorical first-priority claim.

\section{A PL embedding of linear triangle families}
Call a triangle family \emph{linear} if two different triangles share at
most one vertex. Its generated complex contains all their faces.
The following verifies the needed special case of the known PL
hypergraph mechanism, without any separate novelty assertion.

\begin{lemma}\label{lem:PL}
The first barycentric subdivision of the complex generated by a finite
linear triangle family embeds linearly in $\R^3$.
\end{lemma}
\begin{proof}
Form the incidence graph with one node for each original vertex and one
node $f$ for each triangle. Place all nodes on the three-dimensional
moment curve at distinct parameters. Every four nodes are affinely
independent, so this is a straight embedding of the graph.
For a triangle $\{v_1,v_2,v_3\}$ put $a_i=v_i-f$.
Assign its face-barycenter vertex the position $f$, and assign the
edge-barycenter vertex for $\{v_i,v_j\}$ the position
\[
w_{ij}=f+\varepsilon(a_i+a_j),\qquad 0<\varepsilon<\tfrac12.
\]
An edge belongs to just one triangle, so these assignments are globally
consistent. These are positions of the \emph{abstract} subdivision
vertices, not required arithmetic barycenters.

The six subdivision triangles, for $i\ne j$, are
\[
\conv(f,v_i,w_{ij})=\conv(0,e_i,\varepsilon(e_i+e_j))
\]
in the basis $a_1,a_2,a_3$, with origin $f$.
The two in plane $(i,j)$ meet exactly in
$[0,\varepsilon(e_i+e_j)]$. Two different coordinate planes meet on
their common axis; a triangle meets that axis in $[0,e_i]$ if it
contains $e_i$, and otherwise just in $0$.
These are exactly the common subdivision faces, and all triangles
are nondegenerate.

It remains to separate different original triangles. Write $L$ and
$\ell>0$ for the largest and smallest incidence-edge lengths.
Every point in $\conv(f,v_i,w_{ij})$ has the form
\[
x=f+(\lambda+\mu\varepsilon)a_i+\mu\varepsilon a_j,\quad
\lambda,\mu\ge0,\quad\lambda+\mu\le1,
\]
and is within $\varepsilon L$ of the segment $[f,v_i]$.
Choose pairwise disjoint radius-$r$ balls about original vertices,
avoiding facet nodes, such that every nonincident graph edge has
distance greater than $3r$ from the corresponding vertex.
Outside the open radius-$r/2$ balls, the distinct incidence stars
are disjoint compact sets, with a positive minimum separation
$\delta$. Choose $\varepsilon L<\min(r/2,\delta/3)$.
This separates different facet images outside the radius-$r$ balls.

Inside the ball about $v_i$, only triangles containing $v_i$ can
occur. For such a point,
\[
x-v_i=\alpha(f-v_i)+\mu\varepsilon a_j,\qquad
\alpha=1-\lambda-\mu\varepsilon\ge\mu(1-\varepsilon).
\]
Except at $v_i$, its direction is within angle $\arcsin q$ of
$v_i\to f$, where
$q=\varepsilon L/((1-\varepsilon)\ell)<1$.
Distinct facet rays at a shared vertex have a positive minimum
angle $\theta$. Taking also $\arcsin q<\theta/3$ makes their
angular cones disjoint away from the vertex.
All these finitely many conditions admit one positive $\varepsilon$;
conditions for absent pairs are omitted. The local and global
intersection checks give an affine embedding of the whole
subdivision. Incidence cycles cause no additional condition.
\end{proof}

\section{A deletion-resistant affine obstruction}
Put $\mathcal T=\binom{[n]}3$ and $M=|\mathcal T|$.
For a general-position placement $\pi$ of $[n]$ in $\R^4$, let
$W_\pi$ be the unordered pairs of disjoint triangles whose convex
hulls intersect.

Every seven generic points contain a member of $W_\pi$.
For completeness, their homogenized $5\times7$ matrix has a
two-dimensional kernel. On a circle in that kernel, no two
coordinates vanish together, since every five columns are independent.
Along a half-turn the number of positive coordinates changes by one
at each crossing, from $k$ to $7-k$. It therefore crosses between
three and four positives. At that crossing the relation has one zero,
three positive and three negative coefficients, giving the two
intersecting triples. Double counting seven-sets yields
\[
|W_\pi|\ge\frac{\binom n7}{n-6}=\frac{\binom n6}{7}.
\]
A particular triangle belongs to at most $M$ such pairs.

Let $n$ be a sufficiently large integer fourth power, and set
$p=n^{-3/2}$ and $m=n^{5/4}$. Sample each member of $\mathcal T$
independently with probability $p$, obtaining $\mathcal H$.
Fix both a placement order type and a set $F\subseteq\mathcal T$,
$|F|\le m$. Deleting $F$ removes at most $mM$ forbidden pairs.
For $n\ge12$ and $n^{7/4}\ge107520$, the remaining family $W_\pi(F)$
satisfies
\[
|W_\pi(F)|\ge c n^6,\qquad c=\frac1{645120}.
\]
Indeed $\binom n6/7\ge n^6/322560$, and $mM\le n^{17/4}/6$.

Let $Z$ count the pairs in $W_\pi(F)$ both present in $\mathcal H$.
Then $\mu=\mathbb EZ\ge cn^3$.
Two distinct pair indicators are dependent only if they share a
\emph{whole triangle}; sharing geometric vertices alone creates no
dependence. For ordered distinct overlapping pairs the joint
expectation is $p^3$, and their number is at most $M^3$.
Writing their sum as $\Delta$, we have
\[
\mu+\Delta\le M^2p^2/2+M^3p^3\le n^{9/2}.
\]
The standard Janson lower-tail inequality, with this ordered
convention (or equivalently the diagonal included in the denominator),
gives \cite{Janson90,FK}
\[
\Pr(Z=0)\le
\exp\!\left(-\frac{\mu^2}{2(\mu+\Delta)}\right)
\le\exp(-a n^{3/2}),\qquad a=c^2/2.
\]

The labeled simple order type determines $W_\pi$: on each six-set,
the signed five-point minors determine the unique Radon partition.
There are at most $n^{20n}$ such order types in $\R^4$
\cite{GP86,Newman23}. There are at most $(m+1)n^{3m}$ deletion sets.
Consequently the probability that some generic placement is
crossing-free after some deletion of at most $m$ triangles is at most
\begin{equation}\label{eq:failure}
n^{20n}(m+1)n^{3m}\exp(-a n^{3/2})=o(1).
\end{equation}
This union bound explicitly includes deletions chosen \emph{after}
the sample is observed.

Generic placements suffice. In any affine embedding of a finite
complex, the images of all pairs of disjoint nonempty faces have a
positive minimum distance. Small vertex perturbations preserve it.
Any degeneracy or excess intersection between overlapping faces
would, by canceling common barycentric coefficients, give an
intersection between disjoint subfaces. Thus a small generic
perturbation remains an embedding. Unused labels can be added at
generic positions. On the complementary event to \eqref{eq:failure},
every complex generated by $\mathcal H\setminus F$, $|F|\le m$,
has no original-linear embedding in $\R^4$.

\section{Cleanup and the exact minima}
Let $B$ count selected triangle pairs sharing an edge. Then
\[
\mathbb EB=\binom n2\binom{n-2}2p^2\le n/4,\qquad
\Pr(B>m)\le\frac1{4n^{1/4}}.
\]
The triangle count $T$ has mean $\lambda=Mp$ and variance at most
$\lambda$, so $\Pr(T<\lambda/2)\le4/\lambda$.
For large $n$, the three failure probabilities have sum less than
one, and $\lambda/2-m>n$.
Choose a sample where none occurs. From each edge-sharing pair
delete one triangle, taking the union of these choices.
This uses at most $B\le m$ deletions and leaves a linear family
$\mathcal H'$ with more than $n$ triangles.
The choices may depend on the sample.

Let $K$ be its generated complex. Lemma~\ref{lem:PL} gives
$e_{\PL}(K)\le3$ and the claimed first-subdivision embedding.
Its one-skeleton has $3|\mathcal H'|>3n$ distinct edges on at most
$n$ vertices, so it is nonplanar. Hence $e_{\PL}(K)=3$.
The uniform obstruction gives $e_{\lin}(K)\ge5$.
Place its original vertices on the five-dimensional moment curve:
every set of at most six is affinely independent. The union of any
two faces uses at most six vertices, so their convex hulls meet
exactly in their common face. This gives $e_{\lin}(K)\le5$ and proves
Theorem~\ref{thm:main}.

\paragraph{Finite parameters and verification.}
One admissible choice is
$n=2^{256}$, $p=2^{-384}$ and $m=2^{320}$.
The logarithm of the prefactor in \eqref{eq:failure} is less than
$2^{331}$, while $a n^{3/2}>2^{343}$, so the uniform failure is
less than $\exp(-2^{342})<1/8$.
Also $\lambda>2^{379}$, $\lambda/2-m>n$, and the other failures are
at most $2^{-66}$ and less than $2^{-377}$.
Their sum is less than one. The accompanying standard-library
verification package checks these exact bounds and independent
finite geometric controls. It does not enumerate the enormous
witness; the universal arguments above establish existence.

\paragraph{AI use and review status.}
AI tools were used extensively in problem solving, proof verification,
literature review and preparation. Independent AI adversarial audits
and exact arithmetic computations were used. This is an unrefereed
preprint; those checks are not human peer review, and no formally
machine-verified proof is claimed.

\begin{thebibliography}{9}
\bibitem{Brehm06}
U. Brehm, Linear versus piecewise linear embeddability of simplicial
complexes, in \emph{Discrete Differential Geometry}, Oberwolfach
Report 12/2006, pp.~701--702.
\url{https://ems.press/content/serial-article-files/46044}.
\bibitem{BS92}
U. Brehm and K. S. Sarkaria, \emph{Linear vs.\ piecewise-linear
embeddability of simplicial complexes}, MPI/92-52 (1992).
\url{https://archive.mpim-bonn.mpg.de/id/eprint/1946/1/preprint_1992_52.pdf}.
\bibitem{FK}
A. Frieze and M. Karo\'nski, \emph{Introduction to Random Graphs},
author-hosted manuscript dated August 7, 2026,
Section~34.6, Theorem~34.13.
\url{https://www.math.cmu.edu/~af1p/BOOK.pdf}.
\bibitem{GP86}
J. E. Goodman and R. Pollack, Upper bounds for configurations and
polytopes in $\R^d$, \emph{Discrete Comput.\ Geom.} 1 (1986), 219--227.
\href{https://doi.org/10.1007/BF02187696}{doi:10.1007/BF02187696}.
\bibitem{Janson90}
S. Janson, Poisson approximation for large deviations,
\emph{Random Structures Algorithms} 1 (1990), 221--229.
\href{https://doi.org/10.1002/rsa.3240010209}{doi:10.1002/rsa.3240010209}.
\bibitem{LN}
S. Lee and E. Nevo, On colorings of hypergraphs embeddable in $\R^d$,
\emph{Discrete Comput.\ Geom.} 76 (2026), 1942--1959.
Relevant Lemma~3.1 already in arXiv:2307.14195v1 (July 26, 2023);
current preprint v3 (March 9, 2026).
\href{https://doi.org/10.1007/s00454-026-00856-4}{doi:10.1007/s00454-026-00856-4}.
\bibitem{Newman23}
A. Newman, \emph{Linear embeddings of random complexes},
arXiv:2212.09576v3 (October 3, 2023).
\url{https://arxiv.org/pdf/2212.09576v3}.
\end{thebibliography}
\end{document}
